\chapter*{Kết luận}
\label{chap:conclusion}

Khóa luận đã trình bày một cách có hệ thống cơ sở phương pháp luận của các giải thuật trượt gradient ngẫu nhiên tăng cường sử dụng phép thu gọn phương sai dự báo, với trọng tâm là thuật toán SVRG. Về mặt lý thuyết, luận văn đã phân tích chi tiết nguyên nhân khiến SGD cổ điển chịu hội tụ dưới tuyến tính (phương sai gradient không triệt tiêu khi tiến đến nghiệm) và từ đó dẫn dắt ý tưởng thu gọn phương sai thông qua gradient tham chiếu toàn phần định kỳ. Thuật toán SVRG với cấu trúc hai vòng lặp lồng nhau đã được mô tả tường minh, cùng với chứng minh chặt chẽ về tính không chệch và phương sai bị chặn của ước lượng gradient hiệu chỉnh. Khóa luận đã thiết lập định lý hội tụ tuyến tính cho SVRG trên lớp hàm lồi mạnh, chỉ ra tổng số lần đánh giá gradient thành phần để đạt sai số $\epsilon$ là $\mathcal{O}\bigl((n + \kappa)\log(1/\epsilon)\bigr)$, cải thiện đáng kể so với GD và SGD. Về mặt thực nghiệm, chúng tôi đã cài đặt và so sánh SVRG với GD, SGD trên các tập dữ liệu MNIST. Kết quả cho thấy SVRG hội tụ nhanh và ổn định, thời gian thực thi nhanh hơn SGD khoảng một bậc và nhanh hơn GD hàng trăm lần, đồng thời ít nhạy cảm với bước học hơn SGD.

Dù đạt được các mục tiêu đề ra, khóa luận vẫn còn một số hạn chế cần được thừa nhận. Toàn bộ phân tích lý thuyết và phần lớn thực nghiệm chỉ giới hạn trong phạm vi bài toán lồi mạnh, trong khi các hàm mục tiêu trong học sâu thường không lồi. Các thực nghiệm mới chỉ dừng ở mức $n \sim 10^5$; với dữ liệu hàng tỉ mẫu, chi phí tính gradient toàn phần định kỳ có thể trở nên đắt đỏ. Khóa luận chủ yếu tập trung vào SVRG cơ bản, chưa cài đặt và so sánh với các thuật toán hậu duệ mạnh hơn như Katyusha \citep{allen2017katyusha} hay các phương pháp giảm phương sai thích nghi (VR‑Adam, VR‑SGD với động lượng). Ngoài ra, việc tinh chỉnh siêu tham số ($m$, $\eta$) còn được thực hiện thủ công qua grid search, chưa có cơ chế tự thích nghi.

Từ những kết quả và hạn chế đã nêu, một số hướng phát triển tiếp theo đầy triển vọng bao gồm: (i) mở rộng SVRG cho học sâu và bài toán không lồi, kết hợp với các kỹ thuật escape saddle point; (ii) kết hợp SVRG với động lượng và tăng tốc Nesterov như Katyusha \citep{allen2017katyusha} để đạt tốc độ hội tụ tối ưu $O\bigl((n + \sqrt{n\kappa})\log(1/\epsilon)\bigr)$; (iii) phát triển các phương pháp giảm phương sai thích nghi (Adam + Variance Reduction) nhằm tự điều chỉnh bước học theo từng tọa độ; (iv) mở rộng nguyên lý thu gọn phương sai sang học trực tuyến và suy luận Bayes (ví dụ SGLD); (v) triển khai các phiên bản SVRG phân tán hoặc tận dụng GPU để xử lý dữ liệu siêu lớn. Nhìn chung, thu gọn phương sai là một nguyên lý mạnh mẽ và còn nhiều tiềm năng phát triển. Khóa luận này chỉ là bước khởi đầu, và chúng tôi hy vọng sẽ tiếp tục khai thác sâu hơn các hướng đi trên trong tương lai.